\begin{definition}
\label{definition-tor-amplitude}
Let $(X, \mathcal{O}_X)$ be a ringed space.
Let $E$ be an object of $D(\mathcal{O}_X)$.
Let $a, b \in \mathbf{Z}$ with $a \leq b$.
\begin{enumerate}
\item We say $E$ has {\it tor-amplitude in $[a, b]$}
if $H^i(E \otimes_{\mathcal{O}_X}^\mathbf{L} \mathcal{F}) = 0$
for all $\mathcal{O}_X$-modules $\mathcal{F}$ and all $i \not \in [a, b]$.
\item We say $E$ has {\it finite tor dimension}
if it has tor-amplitude in $[a, b]$ for some $a, b$.
\item We say $E$ {\it locally has finite tor dimension}
if there exists an open covering $X = \bigcup U_i$ such that
$E|_{U_i}$ has finite tor dimension for all $i$.
\end{enumerate}
An $\mathcal{O}_X$-module $\mathcal{F}$ has {\it tor dimension $\leq d$}
if $\mathcal{F}[0]$ viewed as an object of $D(\mathcal{O}_X)$ has
tor-amplitude in $[-d, 0]$.
\end{definition}